\documentclass{article}

\usepackage{amsmath, amssymb}
\usepackage{graphicx}
\usepackage{geometry}
\usepackage{fancyhdr}
\usepackage{enumitem}
\usepackage[french]{babel}
\usepackage[utf8]{inputenc}
\usepackage[T1]{fontenc}

\pagestyle{fancy}
\fancyhf{}
\lhead{Correction du Contrôle 2 (A)}
\rhead{Première Spécialité Mathématiques}
\cfoot{\thepage}

\begin{document}

\section*{Correction du Contrôle 2 (A)}

\subsection*{Exercice 1 (6 points)}

\textbf{1. Pour chacune des suites suivantes, indiquer si elle est arithmétique ou géométrique, et préciser pourquoi. Justifier si elle ne l'est pas.}

\begin{enumerate}[label=\alph*)]

\item \( u_1 = 1, \quad u_{n+1} = u_n + \dfrac{7}{4} \)

\textbf{Correction :}

La différence entre deux termes consécutifs est constante :

\[ u_{n+1} - u_n = \dfrac{7}{4} \]

Donc, la suite est \textbf{arithmétique} de raison \( r = \dfrac{7}{4} \).

\bigskip

\item \( u_0 = 1, \quad u_{n+1} = \dfrac{3 u_n}{2} \)

\textbf{Correction :}

Le quotient de deux termes consécutifs est constant :

\[ \dfrac{u_{n+1}}{u_n} = \dfrac{3}{2} \]

Donc, la suite est \textbf{géométrique} de raison \( q = \dfrac{3}{2} \).

\bigskip

\item \( u_0 = 1, \quad u_{n+1} = 2n u_n \)

\textbf{Correction :}

La différence \( u_{n+1} - u_n \) dépend de \( n \), donc n'est pas constante. De même, le quotient \( \dfrac{u_{n+1}}{u_n} = 2n \) dépend de \( n \).

Donc, la suite n'est \textbf{ni arithmétique ni géométrique}.

\bigskip

\item \( u_n = 5 - \pi n \)

\textbf{Correction :}

La suite est définie explicitement. Calculons la différence entre deux termes consécutifs :

\[ u_{n+1} - u_n = [5 - \pi(n+1)] - [5 - \pi n] = -\pi \]

La différence est constante, donc la suite est \textbf{arithmétique} de raison \( r = -\pi \).

\bigskip

\item \( u_n = 2{,}3 \times \left( \dfrac{1}{2} \right)^n \)

\textbf{Correction :}

Le quotient de deux termes consécutifs est :

\[ \dfrac{u_{n+1}}{u_n} = \dfrac{2{,}3 \left( \dfrac{1}{2} \right)^{n+1}}{2{,}3 \left( \dfrac{1}{2} \right)^n} = \dfrac{1}{2} \]

Le quotient est constant, donc la suite est \textbf{géométrique} de raison \( q = \dfrac{1}{2} \).

\bigskip

\item \( u_n = n^2 - 10n \)

\textbf{Correction :}

La différence \( u_{n+1} - u_n \) et le quotient \( \dfrac{u_{n+1}}{u_n} \) ne sont pas constants, ils dépendent de \( n \).

Donc, la suite n'est \textbf{ni arithmétique ni géométrique}.

\end{enumerate}

\bigskip

\textbf{2. Pour chaque suite, déterminer ses variations (croissante, décroissante) sur \( \mathbb{N} \).}

\begin{enumerate}[label=\alph*)]

\item Suite arithmétique de raison \( r = \dfrac{7}{4} > 0 \).

Donc, la suite est \textbf{strictement croissante}.

\bigskip

\item Suite géométrique de raison \( q = \dfrac{3}{2} > 1 \).

Donc, la suite est \textbf{strictement croissante}.

\bigskip

\item Comme \( u_0 = 1 \), calculons \( u_1 \) :

\[ u_1 = 2 \times 0 \times u_0 = 0 \]

Pour \( n \geq 1 \), \( u_n = 0 \) car \( u_{n+1} = 2n u_n = 2n \times 0 = 0 \).

Donc, la suite \textbf{décroît} de \( u_0 = 1 \) à \( u_1 = 0 \), puis reste \textbf{constante} égale à 0.

\bigskip

\item Suite arithmétique de raison \( r = -\pi < 0 \).

Donc, la suite est \textbf{strictement décroissante}.

\bigskip

\item Suite géométrique de raison \( q = \dfrac{1}{2} \) avec \( 0 < q < 1 \).

Donc, la suite est \textbf{strictement décroissante}.

\bigskip

\item \( u_n = n^2 - 10n \)

Calculons la différence \( u_{n+1} - u_n \):

\[ u_{n+1} - u_n = (n+1)^2 - 10(n+1) - [n^2 - 10n] = 2n - 9 \]

- Pour \( n \leq 4 \), \( u_{n+1} - u_n \leq -1 \), la suite est \textbf{décroissante}.
- Pour \( n = 5 \), \( u_{6} - u_{5} = 2 \times 5 - 9 = 1 > 0 \), la suite commence à \textbf{croître}.
- Pour \( n \geq 5 \), \( u_{n+1} - u_n \geq 1 \), la suite est \textbf{croissante}.

Donc, la suite est \textbf{décroissante} pour \( n \leq 4 \), atteint un minimum à \( n = 5 \), puis est \textbf{croissante} pour \( n \geq 5 \).

\end{enumerate}

\bigskip

\textbf{3. Lorsque c'est possible, exprimer chaque suite de manière explicite et calculer son 100\ieme{} terme.}

\begin{enumerate}[label=\alph*)]

\item Suite arithmétique de premier terme \( u_1 = 1 \) et raison \( r = \dfrac{7}{4} \).

La formule explicite est :

\[ u_n = u_1 + (n - 1) r = 1 + (n - 1) \times \dfrac{7}{4} \]

Calcul du 100\ieme{} terme :

\[ u_{100} = 1 + 99 \times \dfrac{7}{4} = 1 + \dfrac{693}{4} = \dfrac{4}{4} + \dfrac{693}{4} = \dfrac{697}{4} \]

\bigskip

\item Suite géométrique de premier terme \( u_0 = 1 \) et raison \( q = \dfrac{3}{2} \).

La formule explicite est :

\[ u_n = u_0 \times q^n = \left( \dfrac{3}{2} \right)^n \]

Calcul du 100\ieme{} terme :

\[ u_{100} = \left( \dfrac{3}{2} \right)^{100} \]

\bigskip

\item La suite n'est ni arithmétique ni géométrique, et il n'est pas possible de l'exprimer simplement de manière explicite.

Cependant, comme \( u_n = 0 \) pour \( n \geq 1 \), on a :

\[ u_{100} = 0 \]

\bigskip

\item La suite est déjà définie explicitement :

\[ u_n = 5 - \pi n \]

Calcul du 100\ieme{} terme :

\[ u_{100} = 5 - 100 \pi \]

\bigskip

\item La suite est déjà définie explicitement :

\[ u_n = 2{,}3 \times \left( \dfrac{1}{2} \right)^n \]

Calcul du 100\ieme{} terme :

\[ u_{100} = 2{,}3 \times \left( \dfrac{1}{2} \right)^{100} \]

\bigskip

\item La suite est déjà définie explicitement :

\[ u_n = n^2 - 10n \]

Calcul du 100\ieme{} terme :

\[ u_{100} = 100^2 - 10 \times 100 = 10\,000 - 1\,000 = 9\,000 \]

\end{enumerate}

\bigskip

\subsection*{Exercice 2 (4 points)}

\textbf{1. Calculer les sommes suivantes :}

\begin{enumerate}[label=\alph*)]

\item \( S = 3 + 7 + 11 + \dots + 487 \)

\textbf{Correction :}

Il s'agit d'une suite arithmétique de premier terme \( u_1 = 3 \) et raison \( r = 4 \).

Trouvons le nombre de termes \( n \):

\[ u_n = u_1 + (n - 1) r \]

\[ 487 = 3 + (n - 1) \times 4 \]

\[ 487 - 3 = 4 (n - 1) \]

\[ 484 = 4 (n - 1) \]

\[ n - 1 = \dfrac{484}{4} = 121 \]

\[ n = 122 \]

La somme des termes est donnée par :

\[ S = \dfrac{n (u_1 + u_n)}{2} \]

\[ S = \dfrac{122 (3 + 487)}{2} = \dfrac{122 \times 490}{2} = 61 \times 490 = 29\,890 \]

\bigskip

\item \( S = 1 + \dfrac{1}{4} + \dfrac{1}{16} + \dots + \dfrac{1}{16\,384} \)

\textbf{Correction :}

Il s'agit d'une suite géométrique de premier terme \( u_0 = 1 \) et raison \( q = \dfrac{1}{4} \).

Le terme général est :

\[ u_n = u_0 \times q^n = \left( \dfrac{1}{4} \right)^n \]

Trouvons le nombre de termes \( n \) tel que :

\[ u_n = \dfrac{1}{16\,384} \]

Sachant que \( 4^{7} = 16\,384 \), donc \( \left( \dfrac{1}{4} \right)^7 = \dfrac{1}{16\,384} \). Ainsi, \( n = 7 \).

La somme des termes est donnée par :

\[ S = u_0 \times \dfrac{1 - q^{n+1}}{1 - q} \]

\[ S = 1 \times \dfrac{1 - \left( \dfrac{1}{4} \right)^{8}}{1 - \dfrac{1}{4}} = \dfrac{1 - \dfrac{1}{4^8}}{\dfrac{3}{4}} = \dfrac{4}{3} \left( 1 - \dfrac{1}{65\,536} \right) \]

\[ S \approx \dfrac{4}{3} \left( 1 - 0 \right) = \dfrac{4}{3} \]

(Note : \( \dfrac{1}{65\,536} \) est très petit et négligeable.)

\bigskip

\item \( S = \dfrac{2}{3} + \dfrac{2}{9} + \dfrac{2}{27} + \dots + \dfrac{2}{2\,187} \)

\textbf{Correction :}

Il s'agit d'une suite géométrique de premier terme \( u_1 = \dfrac{2}{3} \) et raison \( q = \dfrac{1}{3} \).

Le terme général est :

\[ u_n = u_1 \times q^{n - 1} \]

Trouvons le nombre de termes \( n \) tel que :

\[ u_n = \dfrac{2}{2\,187} = \dfrac{2}{3} \times \dfrac{1}{729} \]

Sachant que \( 3^6 = 729 \), donc \( \left( \dfrac{1}{3} \right)^{6} = \dfrac{1}{729} \). Ainsi, \( n - 1 = 6 \), donc \( n = 7 \).

La somme des termes est donnée par :

\[ S = u_1 \times \dfrac{1 - q^{n}}{1 - q} \]

\[ S = \dfrac{2}{3} \times \dfrac{1 - \left( \dfrac{1}{3} \right)^{7}}{1 - \dfrac{1}{3}} = \dfrac{2}{3} \times \dfrac{1 - \dfrac{1}{2\,187}}{\dfrac{2}{3}} \]

\[ S = \left( 1 - \dfrac{1}{2\,187} \right) \]

\[ S \approx 1 \]

(Note : \( \dfrac{1}{2\,187} \) est très petit et négligeable.)

\end{enumerate}

\bigskip

\textbf{2. Exprimer la somme \( \displaystyle \sum_{k=2}^{10} (5k + 1) \) sous la forme \( a + b + \dots + c \).}

\textbf{Correction :}

Calculons les termes pour \( k \) de 2 à 10 :

\begin{align*}
k = 2 &\Rightarrow 5 \times 2 + 1 = 11 \\
k = 3 &\Rightarrow 5 \times 3 + 1 = 16 \\
k = 4 &\Rightarrow 5 \times 4 + 1 = 21 \\
k = 5 &\Rightarrow 5 \times 5 + 1 = 26 \\
k = 6 &\Rightarrow 5 \times 6 + 1 = 31 \\
k = 7 &\Rightarrow 5 \times 7 + 1 = 36 \\
k = 8 &\Rightarrow 5 \times 8 + 1 = 41 \\
k = 9 &\Rightarrow 5 \times 9 + 1 = 46 \\
k = 10 &\Rightarrow 5 \times 10 + 1 = 51 \\
\end{align*}

Donc, la somme est :

\[ 11 + 16 + 21 + 26 + 31 + 36 + 41 + 46 + 51 \]

\bigskip

\subsection*{Exercice 3 (4 points)}

\textbf{1. Dans chaque cas, on donne deux termes d'une suite arithmétique ou géométrique \( (u_n) \). Déterminer la raison et le premier terme \( u_0 \) de cette suite.}

\begin{enumerate}[label=\alph*)]

\item Suite arithmétique, \( u_7 = 31 \) et \( u_{20} = 45 \)

\textbf{Correction :}

La formule générale d'une suite arithmétique est :

\[ u_n = u_0 + n r \]

On a :

\begin{align*}
u_7 &= u_0 + 7 r = 31 \\
u_{20} &= u_0 + 20 r = 45
\end{align*}

En soustrayant les deux équations :

\begin{alignat}{2}
    && (u_{20} - u_7) &= 45 - 31\\
    \iff && (20 r - 7 r) &= 14\\
    \iff && 13r &= 14
\end{alignat}

Donc :

\[ r = \dfrac{14}{13} \]

Ensuite, on calcule \( u_0 \) :

\[ u_0 = u_7 - 7 r = 31 - 7 \times \dfrac{14}{13} = 31 - \dfrac{98}{13} = \dfrac{403}{13} - \dfrac{98}{13} = \dfrac{305}{13} \]

\bigskip

\item Suite géométrique, \( u_3 = 27 \) et \( u_{6} = 729 \)

\textbf{Correction :}

La formule générale d'une suite géométrique est :

\[ u_n = u_0 \times q^n \]

On a :

\begin{align*}
u_3 &= u_0 \times q^3 = 27 \\
u_6 &= u_0 \times q^6 = 729
\end{align*}

En divisant les deux équations :

\[ \dfrac{u_6}{u_3} = \dfrac{u_0 q^6}{u_0 q^3} = q^3 = \dfrac{729}{27} = 27 \]

Donc :

\[ q^3 = 27 \]

\[ q = \sqrt[3]{27} = 3 \]

Ensuite, on calcule \( u_0 \) :

\[ u_3 = u_0 \times q^3 \]

\[ u_0 = \dfrac{u_3}{q^3} = \dfrac{27}{3^3} = \dfrac{27}{27} = 1 \]

\end{enumerate}

\bigskip

\subsection*{Exercice 4 (6 points)}

\textbf{1. En numérotant les étages de haut en bas à partir de \( 0 \), soit \( u_n \) le nombre de cartes à l'étage \( n \), avec \( u_0 = 3 \).}

\begin{enumerate}[label=\alph*)]

\item \textbf{Donner \( u_1 \) et \( u_2 \).}

\textbf{Correction :}

Chaque étage comporte \( 3 \times (n + 1) \) cartes.

- Pour \( n = 1 \) : \( u_1 = 3 \times (1 + 1) = 3 \times 2 = 6 \)
- Pour \( n = 2 \) : \( u_2 = 3 \times (2 + 1) = 3 \times 3 = 9 \)

\bigskip

\item \textbf{Déterminer l'expression de \( u_n \).}

\textbf{Correction :}

\[ u_n = 3 \times (n + 1) \]

\end{enumerate}

\bigskip

\textbf{2. On cherche à calculer le nombre total de cartes \( s_n \) présentes dans un château ayant \( n \) étages.}

\begin{enumerate}[label=\alph*)]

\item \textbf{Donner \( s_0 \), \( s_1 \) et \( s_2 \).}

\textbf{Correction :}

\begin{itemize}
    \item \( s_0 = u_0 = 3 \)
    \item \( s_1 = u_0 + u_1 = 3 + 6 = 9 \)
    \item \( s_2 = u_0 + u_1 + u_2 = 3 + 6 + 9 = 18 \)
\end{itemize}

\bigskip

\item \textbf{Déterminer l'expression de \( s_n \).}

\textbf{Correction :}

\[ s_n = \sum_{k=0}^{n} u_k = \sum_{k=0}^{n} 3(k + 1) = 3 \sum_{k=0}^{n} (k + 1) \]

\[ \sum_{k=0}^{n} (k + 1) = \sum_{k=1}^{n+1} k = \dfrac{(n + 1)(n + 2)}{2} \]

Donc :

\[ s_n = 3 \times \dfrac{(n + 1)(n + 2)}{2} = \dfrac{3(n + 1)(n + 2)}{2} \]

\end{enumerate}

\bigskip

\textbf{3. On dispose de 360 cartes. Combien d'étages peut-on construire ?}

\textbf{Correction :}

On cherche le plus grand \( n \) tel que \( s_n \leq 360 \).

\[ s_n = \dfrac{3(n + 1)(n + 2)}{2} \leq 360 \]

\[ (n + 1)(n + 2) \leq \dfrac{360 \times 2}{3} = 240 \]

Donc :

\[ (n + 1)(n + 2) \leq 240 \]

En testant des valeurs pour \( n \) :

\begin{itemize}
    \item Pour \( n = 13 \) :

    \[ (13 + 1)(13 + 2) = 14 \times 15 = 210 \leq 240 \]

    \item Pour \( n = 14 \) :

    \[ (14 + 1)(14 + 2) = 15 \times 16 = 240 \]

    \item Pour \( n = 15 \) :

    \[ (15 + 1)(15 + 2) = 16 \times 17 = 272 > 240 \]
\end{itemize}

Donc, le nombre maximal d'étages est \( n = 14 \).

\textbf{Réponse :} On peut construire \textbf{14 étages} avec 360 cartes.

\end{document}
